\chapter{Сигналы}
\label{ch:signals}

\section{Что такое сигнал}

\term{Сигнал} --- это физическая величина, которая меняется во времени и
несёт информацию. В электронике чаще всего это \emph{напряжение} на проводе.
Микрофон превращает звук в изменяющееся напряжение, датчик температуры ---
температуру в напряжение, кнопка --- нажатие в напряжение.

Сигналы бывают двух видов.

\begin{description}
  \item[Аналоговый сигнал] может принимать \emph{любое} значение в некотором
    диапазоне: 1,2~В, 1,21~В, 1,215~В\ldots\ Так выглядит звук, сигнал с
    микрофона или датчика.
  \item[Цифровой сигнал] принимает только несколько разрешённых значений. В
    двоичной цифровой технике их два: \textbf{логический 0} (низкое
    напряжение) и \textbf{логическая 1} (высокое напряжение).
\end{description}

\begin{figure}[H]
\centering
\begin{tikzpicture}
\begin{groupplot}[
  group style={group size=2 by 1, horizontal sep=18mm},
  width=0.5\textwidth, height=4.6cm, xmin=0, xmax=10,
  xlabel={Время}, axis lines*=left, xtick=\empty,
  label style={font=\sffamily\small}, tick label style={font=\sffamily\small},
  title style={font=\sffamily\small\bfseries},
]
\nextgroupplot[title={Аналоговый сигнал}, ylabel={$U$, В}, ymin=0, ymax=3.5]
  \addplot[very thick, FpgaBlue, domain=0:10, samples=200] {1.65+1.2*sin(deg(x*0.9))+0.35*sin(deg(x*2.7))};
\nextgroupplot[title={Цифровой сигнал}, ymin=-0.2, ymax=1.4, ytick={0,1}, yticklabels={0,1}, ylabel={логический уровень}]
  \addplot[very thick, FpgaOrange, const plot] coordinates {(0,0) (1,1) (2.5,0) (3,1) (5,1) (5.5,0) (7,1) (8,0) (9,1) (10,1)};
\end{groupplot}
\end{tikzpicture}
\caption{Аналоговый сигнал меняется плавно, цифровой --- скачками между двумя уровнями}
\end{figure}

\subsection{Почему цифровой}

Почему вся современная вычислительная техника цифровая? Главная причина ---
\textbf{устойчивость к помехам}. Любой провод ловит помехи: от соседних
проводов, от моторов, от радиоволн. В аналоговом сигнале помеха искажает
саму информацию, и отделить её от полезного сигнала уже нельзя. В цифровом
сигнале нас интересует только одно: \emph{выше или ниже порога} напряжение.
Небольшая помеха ничего не меняет.

\begin{figure}[H]
\centering
\begin{tikzpicture}
\begin{axis}[
  width=0.9\textwidth, height=5.4cm, xmin=0, xmax=10, ymin=-0.4, ymax=3.8,
  xlabel={Время}, ylabel={$U$, В}, xtick=\empty, ytick={0,0.8,2.0,3.3},
  yticklabels={0, 0{,}8, 2{,}0, 3{,}3},
  axis lines*=left, label style={font=\sffamily\small}, tick label style={font=\sffamily\small},
  clip=false,
]
  \fill[green!12] (axis cs:0,2.0) rectangle (axis cs:10,3.8);
  \fill[red!10] (axis cs:0,0.8) rectangle (axis cs:10,2.0);
  \fill[blue!8] (axis cs:0,-0.4) rectangle (axis cs:10,0.8);
  \addplot[thick, FpgaBlue, samples=400, domain=0:10]
    {(x>1 && x<4 || x>6 && x<8.5 ? 3.1 : 0.2) + 0.25*sin(deg(x*37)) + 0.15*sin(deg(x*91))};
  \node[font=\sffamily\small, anchor=west] at (axis cs:10.1,2.9) {«1»};
  \node[font=\sffamily\small, anchor=west, text=FpgaRed] at (axis cs:10.1,1.4) {запрещено};
  \node[font=\sffamily\small, anchor=west] at (axis cs:10.1,0.2) {«0»};
\end{axis}
\end{tikzpicture}
\caption{Сигнал с помехами. Пока напряжение не заходит в запрещённую зону,
приёмник безошибочно отличает 0 от 1 (пороги --- для стандарта LVCMOS 3,3~В)}
\label{fig:levels}
\end{figure}

\section{Логические уровни}

Какое напряжение считать нулём, а какое --- единицей, определяет
\term{стандарт логических уровней}. ПЛИС на плате Tang Nano 20K использует
стандарт \textbf{LVCMOS33} (питание 3,3~В):

\begin{table}[H]
\centering
\caption{Уровни LVCMOS 3,3 В}
\begin{tabular}{l l l}
\toprule
\textbf{Параметр} & \textbf{Значение} & \textbf{Смысл} \\
\midrule
$U_\text{IL,max}$ & 0,8 В & всё, что ниже, вход считает нулём \\
$U_\text{IH,min}$ & 2,0 В & всё, что выше, вход считает единицей \\
$U_\text{OL,max}$ & 0,4 В & выход выдаёт «0» не выше этого напряжения \\
$U_\text{OH,min}$ & 2,4 В & выход выдаёт «1» не ниже этого напряжения \\
\bottomrule
\end{tabular}
\end{table}

\begin{figure}[H]
\centering
\begin{tikzpicture}[yscale=1.3]
  \fill[blue!10] (0,0) rectangle (2.4,0.4);   \fill[green!15] (0,2.4) rectangle (2.4,3.3);
  \fill[gray!15] (0,0.4) rectangle (2.4,2.4);
  \fill[blue!10] (5,0) rectangle (7.4,0.8);   \fill[green!15] (5,2.0) rectangle (7.4,3.3);
  \fill[red!12] (5,0.8) rectangle (7.4,2.0);
  \draw[thick] (0,0) rectangle (2.4,3.3); \draw[thick] (5,0) rectangle (7.4,3.3);
  \node[font=\sffamily\small\bfseries] at (1.2,3.65) {Выход};
  \node[font=\sffamily\small\bfseries] at (6.2,3.65) {Вход};
  \node[font=\sffamily\small] at (1.2,2.85) {«1»}; \node[font=\sffamily\small] at (1.2,0.2) {«0»};
  \node[font=\sffamily\small] at (6.2,2.65) {«1»}; \node[font=\sffamily\small] at (6.2,0.4) {«0»};
  \node[font=\sffamily\scriptsize, text=FpgaRed] at (6.2,1.4) {неопределённость};
  \foreach \v/\l in {0.4/0{,}4, 2.4/2{,}4, 3.3/3{,}3} \node[font=\sffamily\scriptsize, anchor=east] at (0,\v) {\l~В};
  \foreach \v/\l in {0.8/0{,}8, 2.0/2{,}0} \node[font=\sffamily\scriptsize, anchor=west] at (7.4,\v) {\l~В};
  \draw[FpgaOrange, <->, thick] (3.7,2.0) -- (3.7,2.4) node[midway, right, font=\sffamily\scriptsize, align=left] {запас\\помехо-\\устойчивости};
  \draw[FpgaOrange, <->, thick] (3.7,0.4) -- (3.7,0.8);
  \draw[dashed, FpgaGray] (2.4,2.4) -- (5,2.4); \draw[dashed, FpgaGray] (2.4,2.0) -- (5,2.0);
  \draw[dashed, FpgaGray] (2.4,0.4) -- (5,0.4); \draw[dashed, FpgaGray] (2.4,0.8) -- (5,0.8);
\end{tikzpicture}
\caption{Выход гарантирует уровни «с запасом», поэтому помеха в 0,4 В не
приводит к ошибке}
\end{figure}

\begin{infobox}[Активный высокий и активный низкий уровень]
Обычно «действие» происходит при единице: 1 --- светодиод горит, 1 --- мотор
включён. Это \term{активный высокий уровень}. Но часто бывает наоборот: сигнал
сброса \emph{срабатывает при нуле}. Такие сигналы называют сигналами с
\term{активным низким уровнем} и помечают чертой сверху ($\overline{\text{RESET}}$),
решёткой (RESET\#) или окончанием \code{\_n} (\code{rst\_n}). На плате Tang
Nano 20K светодиоды именно такие --- они горят от нуля.
\end{infobox}

\section{Двоичная система счисления}

Один цифровой провод передаёт один \term{бит} --- 0 или 1. Чтобы передать
число больше единицы, берут несколько проводов (битов) и используют
\term{двоичную систему счисления}. Каждый разряд имеет вес --- степень
двойки:

\begin{figure}[H]
\centering
\begin{tikzpicture}
  \foreach \i/\b in {0/1,1/0,2/1,3/1} {
    \pgfmathtruncatemacro{\p}{3-\i}
    \pgfmathtruncatemacro{\w}{2^(3-\i)}
    \node[draw=FpgaBlue, fill=FpgaLight, thick, minimum size=11mm, font=\ttfamily\Large] (b\i) at (\i*1.6,0) {\b};
    \node[font=\sffamily\small, above=1mm of b\i] {$2^{\p}=\w$};
    \node[font=\sffamily\small, below=1mm of b\i] {$\b\cdot\w$};
  }
  \node[font=\sffamily\small, anchor=west] at (6.2,-0.85) {$= 8+0+2+1 = 11$};
  \node[font=\sffamily\small, anchor=west] at (6.2,0) {$1011_2 = 11_{10}$};
\end{tikzpicture}
\caption{Перевод двоичного числа в десятичное}
\end{figure}

С помощью $n$ битов можно записать $2^n$ разных чисел: от $0$ до $2^n-1$.
Четыре бита дают 16 комбинаций, восемь (\term{байт}) --- 256. Длинные
двоичные числа неудобно читать, поэтому их записывают в
\term{шестнадцатеричной} системе: каждые 4 бита заменяются одной цифрой от 0
до F.

\begin{table}[H]
\centering
\caption{Числа от 0 до 15 в трёх системах счисления}
\small
\begin{tabular}{c c c @{\hspace{12mm}} c c c}
\toprule
\textbf{Дес.} & \textbf{Двоич.} & \textbf{Шестн.} & \textbf{Дес.} & \textbf{Двоич.} & \textbf{Шестн.} \\
\midrule
0 & \code{0000} & 0 & 8  & \code{1000} & 8 \\
1 & \code{0001} & 1 & 9  & \code{1001} & 9 \\
2 & \code{0010} & 2 & 10 & \code{1010} & A \\
3 & \code{0011} & 3 & 11 & \code{1011} & B \\
4 & \code{0100} & 4 & 12 & \code{1100} & C \\
5 & \code{0101} & 5 & 13 & \code{1101} & D \\
6 & \code{0110} & 6 & 14 & \code{1110} & E \\
7 & \code{0111} & 7 & 15 & \code{1111} & F \\
\bottomrule
\end{tabular}
\end{table}

\begin{figure}[H]
\centering
\begin{tikzpicture}
\begin{axis}[
  width=0.75\textwidth, height=5cm, ybar, bar width=6mm,
  xlabel={Число битов $n$}, ylabel={Число комбинаций $2^n$},
  xtick={1,...,8}, ymin=0, ymax=300,
  nodes near coords, nodes near coords style={font=\sffamily\scriptsize},
  label style={font=\sffamily\small}, tick label style={font=\sffamily\small},
  axis lines*=left, grid=major, grid style={FpgaGray!20},
]
  \addplot[fill=FpgaBlue!70, draw=FpgaBlue] coordinates {(1,2) (2,4) (3,8) (4,16) (5,32) (6,64) (7,128) (8,256)};
\end{axis}
\end{tikzpicture}
\caption{Каждый новый бит удваивает число возможных значений}
\end{figure}

\section{Временные диаграммы}

Чтобы показать, как сигналы меняются во времени, рисуют
\term{временные диаграммы}: по горизонтали --- время, по вертикали --- уровень.
Переход из 0 в 1 называют \term{передним фронтом} (или просто фронтом),
переход из 1 в 0 --- \term{задним фронтом} (или спадом).

\begin{figure}[H]
\centering
\begin{tikzpicture}[x=1cm, y=1cm]
  \draw[very thick, FpgaBlue] (0,0) -- (2,0) -- (2.2,1.2) -- (5,1.2) -- (5.2,0) -- (8,0);
  \draw[-{Stealth}, FpgaGray] (0,-0.5) -- (8.5,-0.5) node[right, font=\sffamily\small] {$t$};
  \node[font=\sffamily\small, anchor=east] at (0,0) {0};
  \node[font=\sffamily\small, anchor=east] at (0,1.2) {1};
  \draw[FpgaOrange, thick, -{Stealth}] (1.3,1.9) node[above, font=\sffamily\small] {передний фронт ($0\to1$)} -- (2.05,0.8);
  \draw[FpgaOrange, thick, -{Stealth}] (6.3,1.9) node[above, font=\sffamily\small] {задний фронт, спад ($1\to0$)} -- (5.15,0.8);
  \draw[decorate, decoration={brace, amplitude=5pt, mirror}] (2.2,-0.1) -- (5,-0.1) node[midway, below=5pt, font=\sffamily\scriptsize] {длительность импульса};
\end{tikzpicture}
\caption{Импульс на временной диаграмме}
\end{figure}

Если на диаграмме несколько проводов объединены в число, их рисуют «шиной»:
две линии, между которыми написано значение, а крестик обозначает момент
смены значения.

\begin{figure}[H]
\centering
\begin{tikztimingtable}[timing/slope=0.1, xscale=2.2, yscale=1.4, timing/font=\sffamily\small, timing/rowdist=3.2ex]
  a        & 2L 4H 2L 2H \\
  b        & 4L 4H 2L \\
  data[3:0] & 2D{0} 2D{5} 2D{A} 2D{F} 2D{3} \\
  z (отключён) & 3Z 4H 3Z \\
\end{tikztimingtable}
\caption{Одиночные сигналы, шина и сигнал в третьем (отключённом) состоянии Z}
\end{figure}

\section{Тактовый сигнал}

Особую роль в цифровой технике играет \term{тактовый сигнал} (clock, \code{clk}) ---
последовательность импульсов, которые идут строго через равные промежутки
времени. Он работает как дирижёр: по его фронтам все элементы памяти схемы
одновременно обновляют свои значения.

\begin{figure}[H]
\centering
\begin{tikzpicture}[x=1cm, y=1cm]
  \foreach \i in {0,1,2,3} {
    \draw[very thick, FpgaBlue] (\i*3,0) -- (\i*3,1.2) -- (\i*3+1.5,1.2) -- (\i*3+1.5,0) -- (\i*3+3,0);
    \draw[FpgaOrange, -{Stealth}, thick] (\i*3,1.9) -- (\i*3,1.3);
  }
  \draw[-{Stealth}, FpgaGray] (-0.3,-0.4) -- (12.5,-0.4) node[right, font=\sffamily\small] {$t$};
  \draw[<->, thick] (3,-0.9) -- (6,-0.9) node[midway, below, font=\sffamily\small] {период $T$};
  \draw[<->, thick] (6,1.55) -- (7.5,1.55) node[midway, above, font=\sffamily\small] {$t_\text{и}$};
  \node[font=\sffamily\small, text=FpgaOrange, anchor=west] at (9.3,2.1) {передние фронты};
\end{tikzpicture}
\caption{Тактовый сигнал: период $T$, длительность импульса $t_\text{и}$}
\end{figure}

Основные характеристики тактового сигнала:
\[
  f = \frac{1}{T}, \qquad D = \frac{t_\text{и}}{T}\cdot 100\,\%.
\]
Здесь $f$ --- \term{частота} (сколько периодов в секунду, измеряется в герцах),
$D$ --- \term{коэффициент заполнения} (обычно 50\,\%). На плате Tang Nano 20K
стоит генератор на 27~МГц:
\[
  T = \frac{1}{27\,000\,000\ \text{Гц}} \approx 37 \cdot 10^{-9}\ \text{с} = 37\ \text{нс}.
\]
За одну секунду происходит 27 миллионов тактов. За время одного моргания глаза
(около 0,1~с) --- почти три миллиона.

\begin{table}[H]
\centering
\caption{Приставки для частоты и времени}
\begin{tabular}{l l l l}
\toprule
\textbf{Частота} & \textbf{Период} & \textbf{Пример} \\
\midrule
1 Гц  & 1 с   & мигание светодиода \\
1 кГц = $10^3$ Гц & 1 мс = $10^{-3}$ с & звуковой сигнал «писк» \\
1 МГц = $10^6$ Гц & 1 мкс = $10^{-6}$ с & простые микроконтроллеры \\
27 МГц & 37 нс & генератор Tang Nano 20K \\
1 ГГц = $10^9$ Гц & 1 нс = $10^{-9}$ с & процессоры компьютеров \\
\bottomrule
\end{tabular}
\end{table}

\begin{ideabox}
Цифровой сигнал имеет два уровня: 0 и 1. Несколько битов образуют двоичное
число. Изменения сигналов во времени изображают временными диаграммами, а
тактовый сигнал задаёт ритм, в котором работает вся схема.
\end{ideabox}
